% --- Archon physics formalization source begin ---
% archon:physics
% archon:covers PhyXMiniProblems/problem_phyx_mini_0043.lean
% archon:source-report reports/phyx_mini/problem_phyx_mini_0043.source.json

\chapter{Physics problem phyx\_mini\_0043}
\label{ch:PhyXMiniProblems_problem_phyx_mini_0043}

\paragraph{Problem source.}
\#\# Physical scenario

Converging lenses \textbackslash{}( A \textbackslash{}) and \textbackslash{}( B \textbackslash{}), of focal lengths 8.0\textbackslash{} \textbackslash{}mathrm\{cm\} and 6.0\textbackslash{} \textbackslash{}mathrm\{cm\}, respectively, are placed 36.0\textbackslash{} \textbackslash{}mathrm\{cm\} apart. Both lenses have the same optic axis. An object 8.0\textbackslash{} \textbackslash{}mathrm\{cm\} high is placed 12.0\textbackslash{} \textbackslash{}mathrm\{cm\} to the left of lens \textbackslash{}( A \textbackslash{}).

\#\# Question

Find the magnification of the image produced by the lenses in combination.

\#\# Answer choices

- A: A: -0.856

- B: B: -1.00

- C: C: +0.995

- D: D: -0.814

\#\# Figure caption (auxiliary; use the image as primary evidence)

The image is a diagram of a two-lens system involving Lens A and Lens B.

**Objects and Labels:**
1. **Object O**: Represented by a blue arrow on the left side, placed on the principal axis.
2. **Focal Points**: 
   - \textbackslash{}( F\_1 \textbackslash{}) and \textbackslash{}( F\_2 \textbackslash{}) are the focal points of Lens A.
   - \textbackslash{}( F\_1' \textbackslash{}) and \textbackslash{}( F\_2' \textbackslash{}) are the focal points of Lens B.
3. **Lenses**: 
   - **Lens A** and **Lens B** are depicted as vertical ovals with their principal axes aligned horizontally.

**Rays:**
- Three rays emanate from Object O and pass through Lens A:
  1. Ray 1 (purple) travels parallel to the principal axis to Lens A, passes through the focal point \textbackslash{}( F\_2 \textbackslash{}) on the other side.
  2. Ray 2 (green) passes through the center of Lens A without deviation.
  3. Ray 3 (orange) passes through focal point \textbackslash{}( F\_1 \textbackslash{}) before Lens A and travels parallel to the principal axis afterward.

- Converging at point I to create the first image (I) after Lens A.

- Three rays extend from the image I to pass through Lens B:
  1. Ray 1' (purple) travels parallel to the principal axis to Lens B, passes through the focal point \textbackslash{}( F\_2' \textbackslash{}) on the other side.
  2. Ray 2' (green) passes through the center of Lens B without deviation.
  3. Ray 3' (orange) passes through focal point \textbackslash{}( F\_1' \textbackslash{}) before Lens B and travels parallel to the principal axis afterward.

**Second Image (I')**: Formed by converging rays after passing through Lens B on the right side.

**Measurements:**
- Distances between focal points and lenses, and between images and objects, are marked:
  - 12.0 cm from Lens A to \textbackslash{}( F\_1 \textbackslash{}) and from Lens B to \textbackslash{}( F\_1' \textbackslash{}).
  - 24.0 cm between \textbackslash{}( F\_2 \textbackslash{}) and \textbackslash{}( F\_1' \textbackslash{}).
  - 36.0 cm total distance between Object O and Image I'.
  - Additional smaller segments are marked (8.0 cm, 6.0 cm).

The diagram illustrates the formation of images using light ray

\#\# Dataset metadata

Geometrical Optics; Physical Model Grounding Reasoning, Spatial Relation Reasoning

\paragraph{Recorded answer/context.}
B: B: -1.00

\paragraph{Figure/image path.}
/root/proposal\_for\_physic/science-mango/phyx\_mini\_run/phyx\_data/test\_image/43.png

\paragraph{Formalization target.}
create a compiling Lean file with sorry bodies at `PhyXMiniProblems/problem\_phyx\_mini\_0043.lean`. The Lean declarations must preserve the physical quantities, dimensions or dimensional roles, figure labels, governing-law hypotheses, and final relation expressed by this problem.
Use Mathlib/Physlib names found through LeanExplore where available. If a physics API is missing, introduce faithful local abstractions rather than scalar placeholder aliases.

\begin{theorem}[Physics formalization target]
\label{thm:physics:phyx_mini_0043:target}
\lean{PhyXMiniProblems.ProblemPhyXMini0043.recordedAnswerChoice_ne_overallMagnification}
\uses{def:physics:phyx-mini-0043:phyxminiproblems-problemphyxmini0043-twoconverginglenssetup, def:physics:phyx-mini-0043:phyxminiproblems-problemphyxmini0043-matchessourcereadouts, def:physics:phyx-mini-0043:phyxminiproblems-problemphyxmini0043-matchesfiguregeometry, def:physics:phyx-mini-0043:phyxminiproblems-problemphyxmini0043-matchesraydiagram, def:physics:phyx-mini-0043:phyxminiproblems-problemphyxmini0043-hasphysicalconfiguration, def:physics:phyx-mini-0043:phyxminiproblems-problemphyxmini0043-satisfiesthinlensequationat, def:physics:phyx-mini-0043:phyxminiproblems-problemphyxmini0043-satisfiestransversemagnificationat, def:physics:phyx-mini-0043:phyxminiproblems-problemphyxmini0043-overallmagnification, def:physics:phyx-mini-0043:phyxminiproblems-problemphyxmini0043-displayedmagnification, def:physics:phyx-mini-0043:phyxminiproblems-problemphyxmini0043-recordedanswerchoice}
The assigned autoformalize agent should translate this physics problem into Lean declarations in the covered file, with theorem and lemma proof bodies written as `by sorry`.
\end{theorem}
\begin{proof}
This is an autoformalization task, not a proof task. Produce statements that can later be proved by the physics prover without weakening the physical model.
\end{proof}

% --- Archon declaration topology begin ---
\section{Declaration topology}
\begin{definition}[Signed Length]
  \label{def:physics:phyx-mini-0043:phyxminiproblems-problemphyxmini0043-signedlength}
  \lean{PhyXMiniProblems.ProblemPhyXMini0043.SignedLength}
  A signed physical length, used for coordinates on the optical axis
\end{definition}
\begin{proof}
  This is the declared model definition of Signed Length.
\end{proof}
\begin{definition}[Length Magnitude]
  \label{def:physics:phyx-mini-0043:phyxminiproblems-problemphyxmini0043-lengthmagnitude}
  \lean{PhyXMiniProblems.ProblemPhyXMini0043.LengthMagnitude}
  A nonnegative physical length, used for distances and height magnitudes
\end{definition}
\begin{proof}
  This is the declared model definition of Length Magnitude.
\end{proof}
\begin{definition}[signed Length In Centimeters]
  \label{def:physics:phyx-mini-0043:phyxminiproblems-problemphyxmini0043-signedlengthincentimeters}
  \lean{PhyXMiniProblems.ProblemPhyXMini0043.signedLengthInCentimeters}
  \uses{def:physics:phyx-mini-0043:phyxminiproblems-problemphyxmini0043-signedlength}
  The scalar centimeter readout of a signed physical length
\end{definition}
\begin{proof}
  This is the declared model definition of signed Length In Centimeters.
\end{proof}
\begin{definition}[length Magnitude In Centimeters]
  \label{def:physics:phyx-mini-0043:phyxminiproblems-problemphyxmini0043-lengthmagnitudeincentimeters}
  \lean{PhyXMiniProblems.ProblemPhyXMini0043.lengthMagnitudeInCentimeters}
  \uses{def:physics:phyx-mini-0043:phyxminiproblems-problemphyxmini0043-lengthmagnitude}
  The scalar centimeter readout of a nonnegative physical length
\end{definition}
\begin{proof}
  This is the declared model definition of length Magnitude In Centimeters.
\end{proof}
\begin{definition}[Lens Label]
  \label{def:physics:phyx-mini-0043:phyxminiproblems-problemphyxmini0043-lenslabel}
  \lean{PhyXMiniProblems.ProblemPhyXMini0043.LensLabel}
  The labels of the two lenses, in propagation order
\end{definition}
\begin{proof}
  This is the declared model definition of Lens Label.
\end{proof}
\begin{definition}[Lens Kind]
  \label{def:physics:phyx-mini-0043:phyxminiproblems-problemphyxmini0043-lenskind}
  \lean{PhyXMiniProblems.ProblemPhyXMini0043.LensKind}
  The optical kind stated for each lens
\end{definition}
\begin{proof}
  This is the declared model definition of Lens Kind.
\end{proof}
\begin{definition}[Image Point Label]
  \label{def:physics:phyx-mini-0043:phyxminiproblems-problemphyxmini0043-imagepointlabel}
  \lean{PhyXMiniProblems.ProblemPhyXMini0043.ImagePointLabel}
  The three object/image arrows in the ray diagram
\end{definition}
\begin{proof}
  This is the declared model definition of Image Point Label.
\end{proof}
\begin{definition}[Figure Point Label]
  \label{def:physics:phyx-mini-0043:phyxminiproblems-problemphyxmini0043-figurepointlabel}
  \lean{PhyXMiniProblems.ProblemPhyXMini0043.FigurePointLabel}
  The nine labeled axial points, ordered from left to right in the figure
\end{definition}
\begin{proof}
  This is the declared model definition of Figure Point Label.
\end{proof}
\begin{definition}[Principal Ray Label]
  \label{def:physics:phyx-mini-0043:phyxminiproblems-problemphyxmini0043-principalraylabel}
  \lean{PhyXMiniProblems.ProblemPhyXMini0043.PrincipalRayLabel}
  The six principal rays, with primed labels reserved for lens B
\end{definition}
\begin{proof}
  This is the declared model definition of Principal Ray Label.
\end{proof}
\begin{definition}[Principal Ray Role]
  \label{def:physics:phyx-mini-0043:phyxminiproblems-problemphyxmini0043-principalrayrole}
  \lean{PhyXMiniProblems.ProblemPhyXMini0043.PrincipalRayRole}
  The three standard paraxial constructions shown at each converging lens
\end{definition}
\begin{proof}
  This is the declared model definition of Principal Ray Role.
\end{proof}
\begin{definition}[Image Orientation]
  \label{def:physics:phyx-mini-0043:phyxminiproblems-problemphyxmini0043-imageorientation}
  \lean{PhyXMiniProblems.ProblemPhyXMini0043.ImageOrientation}
  Orientation of a transverse arrow relative to the principal axis
\end{definition}
\begin{proof}
  This is the declared model definition of Image Orientation.
\end{proof}
\begin{definition}[lens]
  \label{def:physics:phyx-mini-0043:phyxminiproblems-problemphyxmini0043-principalraylabel-lens}
  \lean{PhyXMiniProblems.ProblemPhyXMini0043.PrincipalRayLabel.lens}
  \uses{def:physics:phyx-mini-0043:phyxminiproblems-problemphyxmini0043-lenslabel, def:physics:phyx-mini-0043:phyxminiproblems-problemphyxmini0043-principalraylabel}
  The lens traversed by a labeled principal ray
\end{definition}
\begin{proof}
  This is the declared model definition of lens.
\end{proof}
\begin{definition}[role]
  \label{def:physics:phyx-mini-0043:phyxminiproblems-problemphyxmini0043-principalraylabel-role}
  \lean{PhyXMiniProblems.ProblemPhyXMini0043.PrincipalRayLabel.role}
  \uses{def:physics:phyx-mini-0043:phyxminiproblems-problemphyxmini0043-principalraylabel, def:physics:phyx-mini-0043:phyxminiproblems-problemphyxmini0043-principalrayrole}
  The optical construction represented by each numbered ray
\end{definition}
\begin{proof}
  This is the declared model definition of role.
\end{proof}
\begin{definition}[stage Object]
  \label{def:physics:phyx-mini-0043:phyxminiproblems-problemphyxmini0043-stageobject}
  \lean{PhyXMiniProblems.ProblemPhyXMini0043.stageObject}
  \uses{def:physics:phyx-mini-0043:phyxminiproblems-problemphyxmini0043-lenslabel, def:physics:phyx-mini-0043:phyxminiproblems-problemphyxmini0043-imagepointlabel}
  The object point for each of the two imaging stages
\end{definition}
\begin{proof}
  This is the declared model definition of stage Object.
\end{proof}
\begin{definition}[stage Image]
  \label{def:physics:phyx-mini-0043:phyxminiproblems-problemphyxmini0043-stageimage}
  \lean{PhyXMiniProblems.ProblemPhyXMini0043.stageImage}
  \uses{def:physics:phyx-mini-0043:phyxminiproblems-problemphyxmini0043-lenslabel, def:physics:phyx-mini-0043:phyxminiproblems-problemphyxmini0043-imagepointlabel}
  The image point for each of the two imaging stages
\end{definition}
\begin{proof}
  This is the declared model definition of stage Image.
\end{proof}
\begin{definition}[figure Point]
  \label{def:physics:phyx-mini-0043:phyxminiproblems-problemphyxmini0043-imagepointlabel-figurepoint}
  \lean{PhyXMiniProblems.ProblemPhyXMini0043.ImagePointLabel.figurePoint}
  \uses{def:physics:phyx-mini-0043:phyxminiproblems-problemphyxmini0043-imagepointlabel, def:physics:phyx-mini-0043:phyxminiproblems-problemphyxmini0043-figurepointlabel}
  The axial point carrying each object/image-arrow label
\end{definition}
\begin{proof}
  This is the declared model definition of figure Point.
\end{proof}
\begin{definition}[lens Center Point]
  \label{def:physics:phyx-mini-0043:phyxminiproblems-problemphyxmini0043-lenscenterpoint}
  \lean{PhyXMiniProblems.ProblemPhyXMini0043.lensCenterPoint}
  \uses{def:physics:phyx-mini-0043:phyxminiproblems-problemphyxmini0043-lenslabel, def:physics:phyx-mini-0043:phyxminiproblems-problemphyxmini0043-figurepointlabel}
  The center point labeled for each lens
\end{definition}
\begin{proof}
  This is the declared model definition of lens Center Point.
\end{proof}
\begin{definition}[near Focal Point]
  \label{def:physics:phyx-mini-0043:phyxminiproblems-problemphyxmini0043-nearfocalpoint}
  \lean{PhyXMiniProblems.ProblemPhyXMini0043.nearFocalPoint}
  \uses{def:physics:phyx-mini-0043:phyxminiproblems-problemphyxmini0043-lenslabel, def:physics:phyx-mini-0043:phyxminiproblems-problemphyxmini0043-figurepointlabel}
  The incident-side focal point labeled for each lens
\end{definition}
\begin{proof}
  This is the declared model definition of near Focal Point.
\end{proof}
\begin{definition}[far Focal Point]
  \label{def:physics:phyx-mini-0043:phyxminiproblems-problemphyxmini0043-farfocalpoint}
  \lean{PhyXMiniProblems.ProblemPhyXMini0043.farFocalPoint}
  \uses{def:physics:phyx-mini-0043:phyxminiproblems-problemphyxmini0043-lenslabel, def:physics:phyx-mini-0043:phyxminiproblems-problemphyxmini0043-figurepointlabel}
  The outgoing-side focal point labeled for each lens
\end{definition}
\begin{proof}
  This is the declared model definition of far Focal Point.
\end{proof}
\begin{definition}[Thin Lens]
  \label{def:physics:phyx-mini-0043:phyxminiproblems-problemphyxmini0043-thinlens}
  \lean{PhyXMiniProblems.ProblemPhyXMini0043.ThinLens}
  \uses{def:physics:phyx-mini-0043:phyxminiproblems-problemphyxmini0043-lengthmagnitude, def:physics:phyx-mini-0043:phyxminiproblems-problemphyxmini0043-lenslabel, def:physics:phyx-mini-0043:phyxminiproblems-problemphyxmini0043-lenskind}
  A thin lens with its figure label, optical kind, and focal-length magnitude
\end{definition}
\begin{proof}
  This is the declared model definition of Thin Lens.
\end{proof}
\begin{definition}[Two Converging Lens Setup]
  \label{def:physics:phyx-mini-0043:phyxminiproblems-problemphyxmini0043-twoconverginglenssetup}
  \lean{PhyXMiniProblems.ProblemPhyXMini0043.TwoConvergingLensSetup}
  \uses{def:physics:phyx-mini-0043:phyxminiproblems-problemphyxmini0043-signedlength, def:physics:phyx-mini-0043:phyxminiproblems-problemphyxmini0043-lengthmagnitude, def:physics:phyx-mini-0043:phyxminiproblems-problemphyxmini0043-lenslabel, def:physics:phyx-mini-0043:phyxminiproblems-problemphyxmini0043-imagepointlabel, def:physics:phyx-mini-0043:phyxminiproblems-problemphyxmini0043-figurepointlabel, def:physics:phyx-mini-0043:phyxminiproblems-problemphyxmini0043-principalraylabel, def:physics:phyx-mini-0043:phyxminiproblems-problemphyxmini0043-imageorientation, def:physics:phyx-mini-0043:phyxminiproblems-problemphyxmini0043-thinlens}
  This structure records the Two Converging Lens Setup component of the problem-specific physical model
\end{definition}
\begin{proof}
  This is the declared model definition of Two Converging Lens Setup.
\end{proof}
\begin{definition}[signed Height In Centimeters]
  \label{def:physics:phyx-mini-0043:phyxminiproblems-problemphyxmini0043-signedheightincentimeters}
  \lean{PhyXMiniProblems.ProblemPhyXMini0043.signedHeightInCentimeters}
  \uses{def:physics:phyx-mini-0043:phyxminiproblems-problemphyxmini0043-lengthmagnitudeincentimeters, def:physics:phyx-mini-0043:phyxminiproblems-problemphyxmini0043-imagepointlabel, def:physics:phyx-mini-0043:phyxminiproblems-problemphyxmini0043-twoconverginglenssetup}
  Convert an orientation and a height magnitude to a signed height readout
\end{definition}
\begin{proof}
  This is the declared model definition of signed Height In Centimeters.
\end{proof}
\begin{definition}[Matches Source Readouts]
  \label{def:physics:phyx-mini-0043:phyxminiproblems-problemphyxmini0043-matchessourcereadouts}
  \lean{PhyXMiniProblems.ProblemPhyXMini0043.MatchesSourceReadouts}
  \uses{def:physics:phyx-mini-0043:phyxminiproblems-problemphyxmini0043-lengthmagnitudeincentimeters, def:physics:phyx-mini-0043:phyxminiproblems-problemphyxmini0043-twoconverginglenssetup}
  This def records the Matches Source Readouts component of the problem-specific physical model
\end{definition}
\begin{proof}
  This is the declared model definition of Matches Source Readouts.
\end{proof}
\begin{definition}[Matches Figure Geometry]
  \label{def:physics:phyx-mini-0043:phyxminiproblems-problemphyxmini0043-matchesfiguregeometry}
  \lean{PhyXMiniProblems.ProblemPhyXMini0043.MatchesFigureGeometry}
  \uses{def:physics:phyx-mini-0043:phyxminiproblems-problemphyxmini0043-signedlengthincentimeters, def:physics:phyx-mini-0043:phyxminiproblems-problemphyxmini0043-lengthmagnitudeincentimeters, def:physics:phyx-mini-0043:phyxminiproblems-problemphyxmini0043-twoconverginglenssetup}
  This def records the Matches Figure Geometry component of the problem-specific physical model
\end{definition}
\begin{proof}
  This is the declared model definition of Matches Figure Geometry.
\end{proof}
\begin{definition}[Matches Ray Diagram]
  \label{def:physics:phyx-mini-0043:phyxminiproblems-problemphyxmini0043-matchesraydiagram}
  \lean{PhyXMiniProblems.ProblemPhyXMini0043.MatchesRayDiagram}
  \uses{def:physics:phyx-mini-0043:phyxminiproblems-problemphyxmini0043-stageobject, def:physics:phyx-mini-0043:phyxminiproblems-problemphyxmini0043-stageimage, def:physics:phyx-mini-0043:phyxminiproblems-problemphyxmini0043-twoconverginglenssetup}
  This def records the Matches Ray Diagram component of the problem-specific physical model
\end{definition}
\begin{proof}
  This is the declared model definition of Matches Ray Diagram.
\end{proof}
\begin{definition}[Has Physical Configuration]
  \label{def:physics:phyx-mini-0043:phyxminiproblems-problemphyxmini0043-hasphysicalconfiguration}
  \lean{PhyXMiniProblems.ProblemPhyXMini0043.HasPhysicalConfiguration}
  \uses{def:physics:phyx-mini-0043:phyxminiproblems-problemphyxmini0043-lengthmagnitudeincentimeters, def:physics:phyx-mini-0043:phyxminiproblems-problemphyxmini0043-twoconverginglenssetup}
  Nondegeneracy conditions for the real-object, real-image configuration
\end{definition}
\begin{proof}
  This is the declared model definition of Has Physical Configuration.
\end{proof}
\begin{definition}[Satisfies Thin Lens Equation At]
  \label{def:physics:phyx-mini-0043:phyxminiproblems-problemphyxmini0043-satisfiesthinlensequationat}
  \lean{PhyXMiniProblems.ProblemPhyXMini0043.SatisfiesThinLensEquationAt}
  \uses{def:physics:phyx-mini-0043:phyxminiproblems-problemphyxmini0043-lengthmagnitudeincentimeters, def:physics:phyx-mini-0043:phyxminiproblems-problemphyxmini0043-lenslabel, def:physics:phyx-mini-0043:phyxminiproblems-problemphyxmini0043-twoconverginglenssetup}
  This def records the Satisfies Thin Lens Equation At component of the problem-specific physical model
\end{definition}
\begin{proof}
  This is the declared model definition of Satisfies Thin Lens Equation At.
\end{proof}
\begin{definition}[Satisfies Transverse Magnification At]
  \label{def:physics:phyx-mini-0043:phyxminiproblems-problemphyxmini0043-satisfiestransversemagnificationat}
  \lean{PhyXMiniProblems.ProblemPhyXMini0043.SatisfiesTransverseMagnificationAt}
  \uses{def:physics:phyx-mini-0043:phyxminiproblems-problemphyxmini0043-lengthmagnitudeincentimeters, def:physics:phyx-mini-0043:phyxminiproblems-problemphyxmini0043-lenslabel, def:physics:phyx-mini-0043:phyxminiproblems-problemphyxmini0043-stageobject, def:physics:phyx-mini-0043:phyxminiproblems-problemphyxmini0043-stageimage, def:physics:phyx-mini-0043:phyxminiproblems-problemphyxmini0043-twoconverginglenssetup, def:physics:phyx-mini-0043:phyxminiproblems-problemphyxmini0043-signedheightincentimeters}
  This def records the Satisfies Transverse Magnification At component of the problem-specific physical model
\end{definition}
\begin{proof}
  This is the declared model definition of Satisfies Transverse Magnification At.
\end{proof}
\begin{definition}[overall Magnification]
  \label{def:physics:phyx-mini-0043:phyxminiproblems-problemphyxmini0043-overallmagnification}
  \lean{PhyXMiniProblems.ProblemPhyXMini0043.overallMagnification}
  \uses{def:physics:phyx-mini-0043:phyxminiproblems-problemphyxmini0043-twoconverginglenssetup}
  Sequential transverse magnifications multiply for the two-lens system
\end{definition}
\begin{proof}
  This is the declared model definition of overall Magnification.
\end{proof}
\begin{definition}[Answer Choice]
  \label{def:physics:phyx-mini-0043:phyxminiproblems-problemphyxmini0043-answerchoice}
  \lean{PhyXMiniProblems.ProblemPhyXMini0043.AnswerChoice}
  The labels of the four answer choices printed with the problem
\end{definition}
\begin{proof}
  This is the declared model definition of Answer Choice.
\end{proof}
\begin{definition}[displayed Magnification]
  \label{def:physics:phyx-mini-0043:phyxminiproblems-problemphyxmini0043-displayedmagnification}
  \lean{PhyXMiniProblems.ProblemPhyXMini0043.displayedMagnification}
  \uses{def:physics:phyx-mini-0043:phyxminiproblems-problemphyxmini0043-answerchoice}
  The dimensionless magnification printed beside each answer choice
\end{definition}
\begin{proof}
  This is the declared model definition of displayed Magnification.
\end{proof}
\begin{definition}[recorded Answer Choice]
  \label{def:physics:phyx-mini-0043:phyxminiproblems-problemphyxmini0043-recordedanswerchoice}
  \lean{PhyXMiniProblems.ProblemPhyXMini0043.recordedAnswerChoice}
  \uses{def:physics:phyx-mini-0043:phyxminiproblems-problemphyxmini0043-answerchoice}
  The answer label recorded by the dataset
\end{definition}
\begin{proof}
  This is the declared model definition of recorded Answer Choice.
\end{proof}
\begin{lemma}[stage Distances In Centimeters eq]
  \label{lem:physics:phyx-mini-0043:phyxminiproblems-problemphyxmini0043-stagedistancesincentimeters-eq}
  \lean{PhyXMiniProblems.ProblemPhyXMini0043.stageDistancesInCentimeters_eq}
  \uses{def:physics:phyx-mini-0043:phyxminiproblems-problemphyxmini0043-lengthmagnitudeincentimeters, def:physics:phyx-mini-0043:phyxminiproblems-problemphyxmini0043-twoconverginglenssetup, def:physics:phyx-mini-0043:phyxminiproblems-problemphyxmini0043-matchessourcereadouts, def:physics:phyx-mini-0043:phyxminiproblems-problemphyxmini0043-matchesfiguregeometry, def:physics:phyx-mini-0043:phyxminiproblems-problemphyxmini0043-hasphysicalconfiguration, def:physics:phyx-mini-0043:phyxminiproblems-problemphyxmini0043-satisfiesthinlensequationat}
  This lemma records the stage Distances In Centimeters eq component of the problem-specific physical model
\end{lemma}
\begin{proof}
  The conclusion follows from the governing relations and figure readouts named in the statement of stage Distances In Centimeters eq.
\end{proof}
\begin{lemma}[stage Magnifications eq]
  \label{lem:physics:phyx-mini-0043:phyxminiproblems-problemphyxmini0043-stagemagnifications-eq}
  \lean{PhyXMiniProblems.ProblemPhyXMini0043.stageMagnifications_eq}
  \uses{def:physics:phyx-mini-0043:phyxminiproblems-problemphyxmini0043-twoconverginglenssetup, def:physics:phyx-mini-0043:phyxminiproblems-problemphyxmini0043-matchessourcereadouts, def:physics:phyx-mini-0043:phyxminiproblems-problemphyxmini0043-matchesfiguregeometry, def:physics:phyx-mini-0043:phyxminiproblems-problemphyxmini0043-hasphysicalconfiguration, def:physics:phyx-mini-0043:phyxminiproblems-problemphyxmini0043-satisfiesthinlensequationat, def:physics:phyx-mini-0043:phyxminiproblems-problemphyxmini0043-satisfiestransversemagnificationat}
  This lemma records the stage Magnifications eq component of the problem-specific physical model
\end{lemma}
\begin{proof}
  The conclusion follows from the governing relations and figure readouts named in the statement of stage Magnifications eq.
\end{proof}
\begin{theorem}[problem phyx mini 0043]
  \label{thm:physics:phyx-mini-0043:phyxminiproblems-problemphyxmini0043-problem-phyx-mini-0043}
  \lean{PhyXMiniProblems.ProblemPhyXMini0043.problem_phyx_mini_0043}
  \uses{def:physics:phyx-mini-0043:phyxminiproblems-problemphyxmini0043-twoconverginglenssetup, def:physics:phyx-mini-0043:phyxminiproblems-problemphyxmini0043-signedheightincentimeters, def:physics:phyx-mini-0043:phyxminiproblems-problemphyxmini0043-matchessourcereadouts, def:physics:phyx-mini-0043:phyxminiproblems-problemphyxmini0043-matchesfiguregeometry, def:physics:phyx-mini-0043:phyxminiproblems-problemphyxmini0043-matchesraydiagram, def:physics:phyx-mini-0043:phyxminiproblems-problemphyxmini0043-hasphysicalconfiguration, def:physics:phyx-mini-0043:phyxminiproblems-problemphyxmini0043-satisfiesthinlensequationat, def:physics:phyx-mini-0043:phyxminiproblems-problemphyxmini0043-satisfiestransversemagnificationat, def:physics:phyx-mini-0043:phyxminiproblems-problemphyxmini0043-overallmagnification}
  This theorem records the problem phyx mini 0043 component of the problem-specific physical model
\end{theorem}
\begin{proof}
  The conclusion follows from the governing relations and figure readouts named in the statement of problem phyx mini 0043.
\end{proof}
% --- Archon declaration topology end ---
% --- Archon physics formalization source end ---
